\section{2026 年 11 月核验增补}

本节按论文首次提交时间排列。每项方法附有可核验的外部入选依据。

% METHOD: 2609.16443
\paragraph{Neverwhere：从真实场景重建构建视觉跑酷基准}
2026.09~\cite{2609_16443}。Neverwhere 从多视图影像计算相机位姿和多视立体网格，再采样网格初始化高斯。网格承载仿真场景几何，高斯提供视觉观察。构建场景时仍需手工对齐坐标并标注跑酷路线点。

基准把跨障碍、越间隙、斜坡和楼梯拆成逐类任务，检查视觉运动策略成功率和路线进展。这里重建与真实环境之间的差异仍需通过实机试验衡量，基准上的泛化不能直接推出部署表现。

\textbf{入选依据} \href{https://github.com/ziyc/neverwhere}{公开实现}已找到与论文对应的代码文件。

% METHOD: 2609.22809
\paragraph{KIWI：以全景腕部相机采集双手和头部共坐标轨迹}
2026.09~\cite{2609_22809}。KIWI 在两只手上分别布置能同时记录前后方向的相机，后视画面建立共同工作区坐标，前视用于操作观察。将头与双腕的定位统一后，一次采集同时得到高斯场景与双手、头部的六自由度记录。

作者在六种实际操作配置比较纯前视与前后相机的跟踪支持，并用外部标记和独立时钟检查配准、同步。跟踪被支持的帧比例不等同于每帧绝对位姿精度。原文列所有作者均为 Autel US；主题更接近机器人操作采集而非人体头像。

\textbf{入选依据} 论文署名单位包括 Autel（\href{https://arxiv.org/html/2609.22809v1}{论文原文}）。

% METHOD: 2609.35619
\paragraph{CollisionSplatting：用高斯空间不确定性评估避碰代价}
2026.09~\cite{2609_35619}。CollisionSplatting 把场景高斯的空间不确定性纳入距离和避碰代价，利用可调保守系数扩大对潜在障碍的避让。GPU 并行碰撞检查作为模块接入采样规划和模型预测控制，视觉渲染特征可同时用来定义目标。

论文给出仿真中的吞吐量、显存和路径对比，也在真实机械臂和移动机器人上展示了受视觉目标约束的任务。保守度提高通常会带来更多绕行，图中以高斯中心的最近距离做代理指标，实验结果不能当作对新场景无碰撞的保证。

\textbf{入选依据} \href{https://github.com/machines-in-motion/collisionsplatting}{公开实现}已找到与论文对应的代码文件。
